\documentclass[12pt,a4paper]{article}

\usepackage[utf8]{inputenc}
\usepackage{amsmath, amssymb}
\usepackage{graphicx}
\usepackage{hyperref}
\usepackage{geometry}
\usepackage{booktabs}
\usepackage{tcolorbox}
\geometry{a4paper, margin=1in}

\title{\textbf{Viva Material: Ant Colony Optimization for Edge Detection}}
\author{Aditya Kinjawadekar (BT23ECE064)}
\date{\today}

\begin{document}

\maketitle
\tableofcontents
\newpage

% =====================================================================
\section{Edge Detection --- The Basics}
% =====================================================================

\subsection{What is an Edge?}
Think of edges as the outlines in an image. Wherever the brightness of an image changes sharply --- say, where a dark object meets a bright background --- that's an edge. Finding these outlines automatically is one of the most basic and important tasks in image processing.

\subsection{How Do Traditional Methods Find Edges?}
The standard approach is to measure how quickly pixel brightness changes from one pixel to the next. If the change is large, there's probably an edge there.

\begin{itemize}
    \item \textbf{Sobel / Prewitt / Roberts:} These methods slide a small filter (a 3$\times$3 grid of numbers) across the image to estimate the gradient (rate of change) in the horizontal and vertical directions. The edge strength at each pixel is:
    \begin{equation}
        |\nabla I| = \sqrt{G_x^2 + G_y^2}
    \end{equation}
    where $G_x$ is the horizontal gradient and $G_y$ is the vertical gradient.

    \item \textbf{Canny Edge Detector:} A more sophisticated approach that works in stages:
    \begin{enumerate}
        \item Smooth the image with a Gaussian blur to reduce noise.
        \item Compute gradients (same as Sobel).
        \item Thin the edges by keeping only the local maxima (non-maximum suppression).
        \item Apply two thresholds to separate strong edges from weak ones.
        \item Keep weak edges only if they connect to strong edges (hysteresis).
    \end{enumerate}
    Canny works well in practice, but you have to manually pick the right thresholds, and it still processes each pixel independently.
\end{itemize}

\textbf{The core problem with all of these:} every pixel is examined in isolation. There is no ``memory'' of what happened at nearby pixels. This leads to broken, disconnected edge fragments --- especially in noisy images or regions with faint boundaries.

% =====================================================================
\section{Ant Colony Optimization --- The Idea}
% =====================================================================

\subsection{How Real Ants Find the Shortest Path}
When a real ant finds food, it walks back to the colony and leaves a chemical trail called \textbf{pheromone} along the way. Other ants that stumble upon this trail tend to follow it. Ants on shorter paths make more round trips in the same amount of time, so shorter paths accumulate more pheromone. Over time, the colony converges on the shortest route without any single ant ``knowing'' the full picture. This is called \textbf{emergent behavior}.

\subsection{Applying This to Images}
We treat the image as a grid (graph) where each pixel is a node. We release a large number of artificial ants onto this grid. Each ant walks from pixel to pixel, preferring to visit pixels that have:
\begin{enumerate}
    \item High pheromone (other ants have been here before --- it's probably an edge).
    \item High gradient (the pixel itself looks like an edge based on local brightness change).
\end{enumerate}

After enough iterations, the pheromone naturally builds up along the real edges in the image. Pixels that are \textit{not} edges get very little pheromone because ants don't visit them repeatedly.

% =====================================================================
\section{The Math Behind It}
% =====================================================================

There are four key pieces: the heuristic matrix, the transition probability, pheromone evaporation, and pheromone deposit.

\subsection{Heuristic Matrix ($\eta$) --- ``How edge-like does this pixel look?''}
Before any ant starts walking, we compute the Sobel gradient at every pixel. This tells us how strong the intensity change is at that location. We normalize it so the values fall between 0 and 1:
\begin{equation}
    \eta(i,j) = \frac{V(i,j)}{V_{max}}
\end{equation}
$V(i,j)$ is the Sobel gradient magnitude at pixel $(i,j)$. $V_{max}$ is the largest gradient anywhere in the image. A pixel with $\eta$ close to 1 has a very strong gradient --- likely a real edge.

\subsection{Pheromone Matrix ($\tau$) --- ``How many ants have visited this pixel?''}
A matrix the same size as the image, initialized to a small constant (e.g.\ 0.1) everywhere. As ants walk around, they deposit pheromone, so this matrix changes every iteration. It acts as the \textbf{collective memory} of the colony.

\subsection{Transition Probability --- ``Where does an ant go next?''}
An ant sitting at pixel $r$ looks at its 8 neighboring pixels and picks one to move to. The probability of moving to neighbor $s$ is:
\begin{equation}
    P_{r,s} = \frac{[\tau(s)]^\alpha \cdot [\eta(s)]^\beta}{\sum_{u \in \text{neighbors}} [\tau(u)]^\alpha \cdot [\eta(u)]^\beta}
\end{equation}

In plain English: the ant is more likely to move to a neighbor that has \textbf{high pheromone} (other ants liked it) \textit{and} \textbf{high gradient} (it looks like an edge). The denominator just makes sure the probabilities add up to 1.

\textbf{What do $\alpha$ and $\beta$ do?}
\begin{itemize}
    \item $\alpha$ controls how much the ant cares about pheromone (the colony's past experience).
    \item $\beta$ controls how much the ant cares about the gradient (the raw image evidence).
    \item In our project we use $\alpha = 1.0$ and $\beta = 2.0$. Since $\beta > \alpha$, ants prefer strong gradients over pheromone trails. This keeps them anchored to real edges instead of blindly following the crowd.
\end{itemize}

\subsection{Pheromone Update --- ``Evaporate old trails, deposit new ones''}
After all ants finish their walks in a given iteration, two things happen:

\textbf{Step 1: Evaporation.} Old pheromone fades slightly so that bad paths are gradually forgotten:
\begin{equation}
    \tau(i,j) \leftarrow (1 - \rho) \cdot \tau(i,j)
\end{equation}
$\rho$ is the evaporation rate (we use 0.1, meaning 10\% of the pheromone evaporates each round).

\textbf{Step 2: Deposit.} Each ant adds pheromone to the pixels it visited. The amount deposited equals the heuristic value of that pixel:
\begin{equation}
    \tau(i,j) \leftarrow \tau(i,j) + \Delta \tau(i,j)
\end{equation}
This creates a \textbf{positive feedback loop}: strong-gradient pixels get more pheromone, which attracts more ants in the next iteration, which deposits even more pheromone.

% =====================================================================
\section{The Full Algorithm, Step by Step}
% =====================================================================

\begin{enumerate}
    \item Load the image in grayscale.
    \item Compute Sobel gradients to get the heuristic matrix $\eta$.
    \item Initialize the pheromone matrix $\tau$ to a small constant (0.1).
    \item \textbf{For each iteration} (we do 50 in batch mode):
    \begin{enumerate}
        \item Place 8{,}192 ants at random pixel positions.
        \item Each ant takes 10 steps, choosing its next pixel using the transition probability formula above.
        \item After all ants are done, evaporate 10\% of the pheromone and add the new deposits.
    \end{enumerate}
    \item Normalize the final pheromone matrix to the range [0, 255].
    \item Apply a small dilation (3$\times$3 kernel) to thicken the edges for visibility.
    \item Output the result as the edge map.
\end{enumerate}

% =====================================================================
\section{Why ACO is Good for Edge Detection}
% =====================================================================

\begin{tcolorbox}[title=Strengths, colback=blue!5!white, colframe=blue!75!black]
\begin{itemize}
    \item \textbf{Connected edges:} Because ants walk along paths (not pixel by pixel), the resulting edges are more continuous than what Sobel or Canny produce. Gaps get bridged naturally.
    \item \textbf{Noise tolerance:} Random noise doesn't form continuous paths, so ants don't reinforce it. Only real edges get consistent pheromone buildup across many ants and iterations.
    \item \textbf{No training needed:} Unlike deep learning methods, ACO doesn't need labeled datasets or GPUs. It's a pure optimization approach.
    \item \textbf{Tunable:} You can adjust $\alpha$, $\beta$, $\rho$, the number of ants, and the number of iterations to suit different types of images (medical scans, satellite photos, etc.).
\end{itemize}
\end{tcolorbox}

% =====================================================================
\section{Possible Improvements (Good Viva Answers)}
% =====================================================================

If the examiner asks ``How would you make this better?'', here are strong answers:

\subsection{Improvements Within ACO}
\begin{itemize}
    \item \textbf{Adaptive evaporation:} Instead of a fixed $\rho$, decrease it as iterations go on (similar to cooling in simulated annealing). Early iterations explore freely; later ones lock in the best edges.
    \item \textbf{Multi-scale heuristics:} The current version uses a single Sobel filter, which only captures edges at one scale. Using a Gaussian pyramid would let ants detect both fine textures and large-scale boundaries.
    \item \textbf{Hybrid with Canny:} Run Canny first to get a rough edge map, then use ACO only to fill in the gaps. This is much faster because ants don't have to search the entire image.
\end{itemize}

\subsection{Beyond ACO --- Deep Learning Alternatives}
If someone asks about state-of-the-art alternatives:
\begin{itemize}
    \item \textbf{HED (Holistically-Nested Edge Detection):} A CNN that learns edges at multiple scales simultaneously. Produces very clean edges that are immune to texture clutter.
    \item \textbf{DexiNed:} A network that generates crisp, thin edges without needing to be trained on a specific dataset first.
\end{itemize}

\textbf{How to frame this in the viva:} Deep learning is faster (on GPUs) and more accurate, but it needs large datasets and expensive hardware. ACO is slower but fully unsupervised --- it works with just math and no training data. It's a great way to understand how optimization and emergent behavior can solve real problems.

% =====================================================================
\section{Quick Reference: Our Hyperparameters}
% =====================================================================

\begin{center}
\begin{tabular}{llll}
\toprule
\textbf{Parameter} & \textbf{Symbol} & \textbf{Value} & \textbf{What it does} \\
\midrule
Number of ants & $K$ & 8{,}192 & More ants = denser pheromone coverage \\
Iterations & $T$ & 50 & More iterations = better convergence \\
Pheromone weight & $\alpha$ & 1.0 & How much ants follow existing trails \\
Gradient weight & $\beta$ & 2.0 & How much ants follow raw gradients \\
Evaporation rate & $\rho$ & 0.1 & How fast old pheromone fades \\
Steps per ant & $S$ & 10 & How far each ant walks per iteration \\
\bottomrule
\end{tabular}
\end{center}

% =====================================================================
\section{Reference}
% =====================================================================

\begin{enumerate}
    \item ``Image Segmentation Technology Based on Ant Colony Algorithm.'' \textit{Journal of Electrical Systems}, Vol.\ 20 No.\ 9s (2024).\\
    \url{https://journal.esrgroups.org/jes/article/view/4307}
\end{enumerate}

\end{document}
